\section{Moment Generating Functions}

\begin{definition}[Moment Generating Function]
    \label{def:moment-generating-fn}%
    Let \(X\) be a random variable with density \(f\). The \emph{moment generating function} of \(X\) is defined to be
    \[
        m\para{\theta} = \Expt\brac{e^{\theta X}} = \int_{\RR} e^{\theta x} f\para{x} \Diff x
    \]
    whenever this is well-defined, and by convention we let \(m\para{0} = 1\).
\end{definition}

\begin{theorem}[Uniqueness of Moment Generating Function]
    \label{thm:uniqueness-mgf}%
    The moment generating function uniquely determines the distribution of a random variable, provided it is defined for an open interval of values of \(\theta\).
\end{theorem}

\begin{theorem}[Derivative of Moment Generating Function gives the Moment]
    \label{thm:derivative-mgf-moment}%
    Suppose the moment generating function is defined for an open interval of values of \(\theta\). Then
    \[
        m^{\para{r}} \para{0} = \rightbar{\DiffNFrac{m\para{\theta}}{\theta}{r}}_{\theta = 0} = \Expt\brac{X^r}.
    \]
\end{theorem}

We may very easily find the moment generating function of the sum of independent random variables.

\begin{claim}[Moment Generating Function of Sum of Independent Random Variables]
    \label{claim:mgf-sum-independent}%
    If the random variables \(X_1, \ldots, X_n\) are independent, then
    \[
        m\para{\theta} = \Expt\brac{e^{\theta \para{X_1 + \cdots + X_n}}} = \prod_{i = 1}^{n} \Expt\brac{e^{\theta X_i}}.
    \]
\end{claim}

\subsection{Moment Generating Function of Standard Distributions}

\begin{definition}[Gamma Distribution]
    \label{def:gamma-distribution}%
    The \emph{Gamma Distribution} with parameters \(n \in \NN\) and \(\lambda > 0\), denoted by \(\Gamma\para{n, \lambda}\), is given by the density function
    \[
        f\para{x} = \frac{e^{-\lambda x} \cdot \lambda^n \cdot x^{n - 1}}{\para{n - 1}!}
    \]
    for \(x > 0\).
\end{definition}

\begin{proposition}[Gamma Distribution is a Distribution]
    \label{prop:gamma-distribution-integral}%
    The Gamma Distribution is a distribution. In particular, its density function integrates to \(1\).
\end{proposition}

\begin{proof}
    Let \(X \sim \Gamma\para{n, \lambda}\) have density \(f\). Then
    \[
        I_n = \int_{0}^{\plusinfty} f\para{x} \Diff x = \int_{0}^{\plusinfty}  \lambda e^{-\lambda x} \cdot \frac{\lambda^{n - 1} x^{n - 1}}{\para{n - 1}!} \Diff x.
    \]

    Hence,
    \begin{align*}
        I_n &= \int_{0}^{\plusinfty} e^{-\lambda x} \cdot \frac{\para{n - 1} \cdot \lambda^{n - 1} x^{n - 2}}{\para{n - 1}!} \Diff x \\
        &= \int_{0}^{\plusinfty} \frac{e^{-\lambda x} \cdot \lambda^{n - 1} \cdot x^{n - 2}}{\para{n - 2}!} \Diff x \\
        &= I_{n - 1},
    \end{align*}
    so all \(I_n\)s are equal to \(I_1\). We have
    \[
        I_1 = \int_{0}^{\plusinfty} \lambda e^{-\lambda x} \Diff x = 1,
    \]
    so \(f\) is a density.
\end{proof}

\begin{example}[Moment Generating Function of Gamma Distribution]
    \label{ex:gamma-distribution-mgf}%
    We have
    \begin{align*}
        m\para{\theta} = \Expt\brac{e^{\theta X}} &= \int_{0}^{\plusinfty} e^{\theta x} \cdot e^{- \lambda x} \cdot \frac{\lambda^n \cdot x^{n - 1}}{\para{n - 1}!} \Diff x \\
        &= \int_{0}^{\plusinfty} e^{- \para{\lambda - \theta} x} \cdot \frac{\para{\lambda - \theta}^n \cdot x^{n - 1}}{\para{n - 1}!} \Diff x \cdot \frac{\lambda^n}{\para{\lambda - \theta}^n}.
    \end{align*}

    Therefore, if \(\theta < \lambda\), then
    \[
        m \para{\theta} = \para{\frac{\lambda}{\lambda - \theta}}^n
    \]
    for \(\theta < \lambda\).
\end{example}

\begin{example}[Sum of Gamma Random Variables]
    \label{ex:sum-gamma}%
    Let \(X \sim \Gamma\para{n, \lambda}\) and \(Y \sim \Gamma\para{m, \lambda}\), with \(n, m \in \NN\) and \(\lambda > 0\), and \(X \indp Y\).

    We have
    \begin{align*}
        \Expt\brac{e^{\theta \para{X + Y}}} &= \Expt\brac{e^{\theta X}} \cdot \Expt\brac{e^{\theta Y}} \\
        &= \para{\frac{\lambda}{\lambda - \theta}}^n \cdot \para{\frac{\lambda}{\lambda - \theta}}^m \\
        &= \para{\frac{\lambda}{\lambda - \theta}}^{n + m},
    \end{align*}
    and hence \(X + Y \sim \Gamma\para{n + m, \lambda}\).
\end{example}

\begin{example}[Sum of Exponential Random Variables]
    \label{ex:sum-exponential}%
    Let \(X_1, \ldots, X_n\) be independent and identically distributed \(\Exponential\para{\lambda} \sim \Gamma\para{1, \lambda}\) random variables, with \(\lambda > 0\). Then we have
    \[
        X_1 + \cdots + X_n \sim \Gamma\para{n, \lambda}.
    \]
\end{example}

\begin{remark}
    We specified for \(\Gamma\para{n, \lambda}\), that \(n \in \NN\). In fact, we may define \(\Gamma\para{\alpha, \lambda}\) for \(\alpha > 0\), where \(n\) is replaced by \(\alpha\), with the density function
    \[
        f\para{x} = \frac{e^{-\lambda x} \cdot \lambda^{\alpha} \cdot x^{\alpha - 1}}{\Gamma\para{\alpha}}
    \]
    where \(\Gamma\) is the \emph{Gamma Function}
    \[
        \Gamma\para{\alpha} = \int_{0}^{\plusinfty} e^{-x} \cdot x^{\alpha - 1} \Diff x.
    \]
\end{remark}

\begin{definition}[Cauchy Distribution]
    The \emph{Cauchy Distribution} is the continuous random variable defined by the density function
    \[
        f\para{x} = \frac{1}{\pi \para{1 + x^2}}
    \]
    for \(x \in \RR\).
\end{definition}

\begin{remark}
    For any non-zero \(\theta\), we have
    \[
        m\para{\theta} = \int_{\minusinfty}^{\plusinfty} \frac{e^{\theta x}}{\pi \para{1 + x^2}} \Diff x = \plusinfty,
    \]
    while \(m\para{0} = 1\).
\end{remark}

Note that multiples of the Cauchy distribution therefore would have the same moment generating function -- so the condition on the moment generating functions to agree on an \emph{open set} of values is essential.

\begin{example}[Moment Generating Function of Normal Distribution]
    \label{ex:mgf-normal}%
    For a normal distribution \(X \sim \Normal\para{\mu, \sigma^2}\), we have
    \[
        f\para{x} = \frac{1}{\sqrt{2 \pi \sigma^2}} \exp\para{- \frac{\para{x - \mu}^2}{2 \sigma^2}}
    \]
    for \(x \in \RR\). Therefore, we have
    \begin{align*}
        m\para{\theta} &= \Expt\brac{e^{\theta X}} \\
        &= \int_{\minusinfty}^{\infty} e^{\theta x} \cdot \frac{1}{\sqrt{2 \pi \sigma^2}} \exp\para{- \frac{\para{x - \mu}^2}{2 \sigma^2}} \Diff x.
    \end{align*}

    We may simplify the exponent as
    \begin{align*}
        \theta x - \frac{\para{x - \mu}^2}{2 \sigma^2} &= \theta x - \frac{x^2}{2 \sigma^2} + \frac{2 x \mu}{2 \sigma^2} - \frac{\mu^2}{2 \sigma^2} \\
        &= - \frac{x^2}{2 \sigma^2} + 2 \cdot \frac{x}{2 \sigma^2} \para{\mu + \theta \sigma^2} - \frac{\mu^2}{2 \sigma^2} \\
        &= - \frac{x^2}{2 \sigma^2} + 2 \cdot \frac{x}{2 \sigma^2} \para{\mu + \theta \sigma^2} - \frac{\para{mu + \theta \sigma^2}^2}{2 \sigma^2} + \frac{\para{\mu + \theta \sigma^2}^2}{2 \sigma^2} - \frac{\mu^2}{2 \sigma^2} \\
        &= - \frac{1}{2 \sigma^2} \para{x - \para{\mu + \theta \sigma^2}}^2 + \frac{\mu^2}{2 \sigma^2} + \frac{2 \mu \theta \sigma^2}{2 \sigma^2} + \frac{\theta^2 \sigma^2}{2} - \frac{\mu^2}{2 \sigma^2} \\
        &= - \frac{1}{2 \sigma^2} \para{x - \para{\mu + \theta \sigma^2}}^2 + \mu \theta + \frac{\theta^2 \sigma^2}{2}.
    \end{align*}

    Therefore, we have
    \begin{align*}
        m\para{\theta} &= \int_{\minusinfty}^{\plusinfty} \frac{1}{\sqrt{2 \pi \sigma^2}} \exp\para{- \frac{\para{x - \para{\mu + \theta \sigma^2}}^2}{2 \sigma^2}} \Diff x  \cdot \exp \para{\mu \theta + \frac{\theta^2 \sigma^2}{2}} \\
        &= \exp\para{\mu \theta + \frac{\theta^2 \sigma^2}{2}}.
    \end{align*}

    Now, suppose \(X \sim \Normal\para{\mu, \sigma^2}\) and \(Y \sim \Normal\para{\nu, \tau^2}\), and \(X \indp Y\).

    Then, we have
    \begin{align*}
        m\para{\theta} &= \Expt\brac{e^{\theta \para{X + Y}}} \\
        &= \exp\para{\mu \theta + \frac{\theta^2 \sigma^2}{2}} \cdot \exp\para{\nu \theta + \frac{\theta^2 \tau^2}{2}} \\
        &= \exp\para{\para{\mu + \nu} \theta + \frac{\theta^2}{2} \para{\sigma^2 + \tau^2}}.
    \end{align*}

    Therefore, \(X + Y \sim \Normal\para{\mu + \nu, \sigma^2 + \tau^2}\).
\end{example}

\subsection{Multivariate Moment Generating Functions}

\begin{definition}[Multivariate Moment Generating Function]
    \label{def:multivariate-mgf}%
    Let \(\vm{X} = \para{X_1, \ldots, X_n} \in \RR^n\) be an \(n\)-dimensional random variable. Its \emph{moment generating function} is defined to be
    \[
        m\para{\vmb{\theta}} = \Expt\brac{e^{\vmb{\theta}\transpose \vm{X}}} = \Expt\brac{e^{\sum_{i = 1}^{n} \theta_i X_i}}
    \]
    where \(\vmb{\theta} = \para{\theta_1, \ldots, \theta_n} \transpose\).
\end{definition}

\begin{theorem}[Uniqueness of Multivariate Moment Generating Function]
    \label{thm:uniqueness-multivarite-mgf-moments}%
    If the moment generating function is finite for an open set of values of \(\vmb{\theta}\), then it uniquely determines the distribution, and we have
    \[
        \rightbar{\PartialNFrac{m}{\theta_i}{r}}_{\vmb{\theta} = \vm{0}} = \Expt\brac{X_i^r}
    \]
    and
    \[
        \rightbar{\frac{\partial^{r + s} m}{\partial \theta_i^r \partial \theta_j^s}}_{\vmb{\theta} = \vm{0}} = \Expt\brac{X_i^r \cdot X_j^s}.
    \]
\end{theorem}

\begin{claim}[Independent Random Variables give Separable Moment Generating Function]
    \label{claim:independent-rv-separable-mgf}%
    Let \(\vm{X} = \para{X_1, \ldots, X_n} \in \RR^n\) be a random variable. Then
    \[
        m\para{\vmb{\theta}} = \prod_{i = 1}^{n} \Expt\brac{e^{\theta_i X_i}}
    \]
    if and only if \(X_1, \ldots, X_n\) are independent.
\end{claim}

\begin{proof}
    If they are independent, this statement is straightforward.

    If \(m\para{\vmb{\theta}}\) factorises, then by \autoref{thm:uniqueness-multivarite-mgf-moments} (Uniqueness of Moment Generating Functions), we get that this \(\vm{X}\) agrees with the distribution when the \(X_i\)s are independent, so they must be independent.
\end{proof}

\subsection{Multivariate Gaussian Random Variable}

\begin{definition}[Multivariate Gaussian Random Variable]
    \label{def:multivariate-gaussian-rv}%
    A random variable \(X\) with values in \(\RR\) is \emph{Gaussian}, if it could be written as
    \[
        X = \mu + \sigma Z
    \]
    where \(Z \sim \Normal\para{0, 1}\) and \(\mu \in \RR, \sigma \geq 0\). (This includes all normal distributions defined before and the trivial distribution.)

    Let \(\vm{X} = \para{X_1, \ldots, X_n} \in \RR^n\) be a random variable. Then \(\vm{X}\) is \emph{Gaussian}, if for any \(\vm{u} = \para{u_1, \ldots, u_n}\transpose \in \RR^n\), the dot product
    \[
        \vm{u}\transpose \vm{X} = \sum_{i = 1}^{n} u_i X_i
    \]
    is a Gaussian random variable in \(\RR\).
\end{definition}

\begin{claim}[Affine Transformation of Gaussian is Gaussian]
    \label{claim:affine-transformation-preserves-gaussian}%
    Let \(\vm{X} = \para{X_1, \ldots, X_n} \in \RR^n\) be Gaussian. Let \(\vm{A}\) be an \(m \times n\) matrix, and \(\vm{b} \in \RR^m\). Then the result
    \[
        \vm{A} \vm{X} + \vm{b}
    \]
    is also Gaussian.
\end{claim}

\begin{proof}
    Let \(\vm{u} = \para{u_1, \ldots, u_m} \transpose \in \RR^m\) be arbitrary. We want to show \(\vm{u}\transpose \para{\vm{A} \vm{X} + \vm{b}}\) is Gaussian in \(\RR\). We have
    \[
        \vm{u}\transpose \para{\vm{A} \vm{X} + \vm{b}} = \vm{u}\transpose \vm{A} \vm{X} + \vm{u}\transpose \vm{b}.
    \]

    Let \(\vm{v} = \vm{A}\transpose \vm{u}\). Then we have
    \[
        \vm{u}\transpose \para{\vm{A} \vm{X} + \vm{b}} = \vm{v}\transpose \vm{X} + \vm{u}\transpose \vm{b}.
    \]

    \(\vm{v}\transpose \vm{X}\) is Gaussian in \(\RR\), since \(\vm{X}\) is Gaussian. Adding a constant to a Gaussian will still produce a Gaussian in \(\RR\). So
    \[
        \vm{u}\transpose \para{\vm{A} \vm{X} + \vm{b}}
    \]
    is Gaussian in \(\RR\), as desired.
\end{proof}

\begin{proposition}[Moment Generating Function for Multivariate Gaussian]
    \label{prop:mgf-multivariate-gaussian}%
    Let \(\vm{X} = \para{X_1, \ldots, X_n}\) be a multivariate Gaussian random variable. Then its moment generating function, \(m\para{\vmb{\lambda}}\), is given by
    \[
        m\para{\vmb{\lambda}} = \Expt\brac{e^{\vmb{\lambda} \transpose \vm{X}}} = \exp\para{\vmb{\lambda} \transpose \vmb{\mu} + \frac{\vmb{\lambda} \transpose \vm{V} \vmb{\lambda}}{2}},
    \]
    where
    \[
        \vmb{\mu} = \Expt\brac{\vm{X}}, \quad \vm{V} = \Var\para{\vm{X}}.
    \]
\end{proposition}

\begin{proof}
    Note that \(\vmb{\lambda \transpose \vm{X}}\) is a normal distribution with parameters
    \[
        \Normal\para{\vmb{\lambda}\transpose\vmb{\mu}, \vmb{\lambda}\transpose \vm{V} \vmb{\lambda}}
    \]
    by \autoref{prop:expt-var-transform-multivariate-rv}. The result follows.
\end{proof}

Since the moment generating function (on an open set) uniquely characterises the distribution, it is sufficient to know the mean and the variance to characterise the multivariate Gaussian distribution.

\begin{notation}
    We write the Multivariate Gaussian distribution \(\vm{X}\) with mean \(\vmb{\mu} \in \RR^n\), and non-negative definite matrix \(\vm{V}\), as
    \[
        \vm{X} \sim \Normal\para{\vmb{\mu}, \vm{V}}.
    \]
\end{notation}

The next question is, does the multivariate Gaussian always exist, for any specified mean and variance? We start with the simple case. 

\begin{lemma}[Standard Normals produce a Multivariate Gaussian]
    \label{lem:standard-normal-multivariate-gaussian}%
    Let \(Z_1, \ldots, Z_n\) be independent and identically distributed \(\Normal\para{0, 1}\) random variables. Then the multivariate random variable
    \[
        \vm{Z} = \para{Z_1, \ldots, Z_n}
    \]
    is Gaussian. In particular,
    \[
        \vm{Z} \sim \Normal\para{\vm{0}, \identityM_n}.
    \]
\end{lemma}

\begin{proof}
    We find the moment generating function of \(\vm{u} \transpose \vm{Z}\). We have
    \begin{align*}
        m\para{\lambda} &= \Expt\brac{e^{\lambda \vm{u}\transpose \vm{Z}}} \\
        &= \Expt\brac{e^{\lambda \sum_{i = 1}^{n} u_i Z_i}} \\
        &= \prod_{i = 1}^{n} \Expt\brac{e^{\lambda u_i Z_i}} \\
        &= \prod_{i = 1}^{n} \exp\para{\frac{\lambda^2 u_i^2}{2}} \\
        &= \exp\para{\frac{\lambda^2 \abs{\vm{u}}^2}{2}}.
    \end{align*}

    Hence,
    \[
        \vm{u}\transpose \vm{Z} \sim \Normal \para{0, \abs{\vm{u}}^2}
    \]
    is normal, with the desired parameters, as desired.
\end{proof}

\begin{lemma}[Construction of Multivariate Gaussian]
    \label{lem:construction-multivariate-gaussian}%
    Let \(\vmb{\mu} \in \RR^n\), and \(\vm{V}\) be a non-negative definite matrix. Let \(Z_1, \ldots, Z_n\) be identically and independently distributed with \(\Normal\para{0, 1}\). Let \(\vm{Z} = \para{Z_1, \ldots, Z_n}\).

    Then the random variable
    \[
        \vm{X} = \vmb{\mu} + \vmb{\sigma} \vm{Z}
    \]
    is Gaussian, with \(X \sim \Normal\para{\vmb{\mu}, \vm{V}}\), where \(\vmb{\sigma}\) is the square-root of \(\vm{V}\).
\end{lemma}

\begin{proof}
    First, \(\vm{X}\) is Gaussian since it is a linear transformation of a Gaussian vector. We have \(\Expt\brac{\vm{X}} = \vmb{\mu}\). For the variance, we have
    \begin{align*}
        \Var\para{\vm{X}} &= \Expt\brac{\para{\vm{X} - \vmb{\mu}} \cdot \para{\vm{X} - \vmb{\mu}}\transpose} \\
        &= \Expt\brac{\para{\vmb{\sigma} \cdot \vm{Z}} \cdot \para{\vm{Z}\transpose \cdot \vmb{\sigma}\transpose}} \\
        &= \vmb{\sigma} \cdot \Expt\brac{\vm{Z} \cdot \vm{Z}\transpose} \vmb{\sigma}\transpose \\
        &= \vmb{\sigma} \cdot \identityM_n \cdot \vmb{\sigma}\transpose \\
        &= \vmb{\sigma} \cdot \vmb{\sigma} \\
        &= \vm{V}.
    \end{align*}
\end{proof}

We aim to find the density of a multivariate Gaussian next.

\begin{proposition}[Density of Multivariate Gaussian with Positive Definite Variance]
    If \(\vm{V}\) is positive definite, then the density of \(\Normal\para{\vmb{\mu}, \vm{V}}\) is given by
    \[
        f_{\vm{X}}\para{\vm{x}} = \frac{1}{\sqrt{\para{2\pi}^n \det \vm{V}}} \exp\para{- \frac{\para{\vm{x} - \vmb{\mu}}\transpose \vm{V}\inverse \para{\vm{x} - \vmb{\mu}}}{2}}.
    \]
\end{proposition}

\begin{proof}
    We have \(\vm{X} = \vmb{\mu} + \vmb{\sigma} \vm{Z}\), where \(\vm{Z} \sim \Normal\para{\vm{0}, \identityM_n}\).

    Therefore, \(\vm{Z} = \vmb{\sigma}\inverse \para{\vm{X} - \vmb{\mu}}\).

    Hence,
    \[
        f_{\vm{X}} \para{\vm{x}} = f_{\vm{Z}} \para{\vmb{\sigma}\inverse \para{\vm{x} - \vmb{\mu}}} \cdot \abs{\det \vm{J}},
    \]
    where \(\vm{J} = \vmb{\sigma}\inverse\).

    Hence,
    \begin{align*}
        f_{\vm{X}} \para{\vm{x}} &= f_{\vm{Z}} \para{\vmb{\sigma}\inverse \para{\vm{x} - \vmb{\mu}}} \cdot \abs{\det \vm{J}} \\
        &= f_{\vm{Z}} \para{z_1, \cdots, z_n} \cdot \det\para{\vmb{\sigma} \inverse} \\
        &= \frac{1}{\para{2 \pi}^{n / 2}} \exp\para{- \frac{\abs{\vm{z}}^2}{2}} \cdot \det\para{\vmb{\sigma}\inverse} \\
        &= \frac{1}{\para{2 \pi}^{n / 2}} \exp\para{- \frac{\vm{z}\transpose \vm{z}}{2}} \frac{1}{\det \vmb{\sigma}}.
    \end{align*}

    Hence,
    \[
        f_{\vm{X}}\para{\vm{x}} = \frac{1}{\para{2 \pi}^{n / 2} \det \vmb{\sigma}} \exp\para{- \frac{\para{\vm{x} - \vmb{\mu}}\transpose \cdot \para{\vmb{\sigma}\inverse}\transpose \cdot \vmb{\sigma}\inverse \cdot \para{\vm{x} - \vmb{\mu}}}{2}}.
    \]

    We have \(\para{\vmb{\sigma}\inverse}\transpose = \vmb{\sigma}\inverse\), and \(\det\vmb{\sigma} = \sqrt{\det \vm{V}}\), and hence
    \[
        f_{\vm{X}}\para{\vm{x}} = \frac{1}{\sqrt{\para{2\pi}^n \det \vm{V}}} \exp\para{- \frac{\para{\vm{x} - \vmb{\mu}}\transpose \vm{V}\inverse \para{\vm{x} - \vmb{\mu}}}{2}}
    \]
    as desired.
\end{proof}

Now we look at the degenerate case, where one or more of the entries \(\lambda_i = 0\). By an orthogonal change of basis, we can assume that \(\vm{V}\) takes the form of
\[
    \vm{V} = \begin{pmatrix}
        \vm{U} & 0 \\ 0 & 0
    \end{pmatrix}
\]
where \(\vm{U}\) is a positive definite \(m \times m\) matrix, where \(m < n\).

Under this change of basis, we have
\[
    \vmb{\mu} = \begin{pmatrix}
        \vmb{\lambda} \\ \vmb{\nu}
    \end{pmatrix}
\]
for \(\vmb{\lambda} \in \RR^m, \vmb{\nu} \in \RR^{n - m}\). Hence, we can write the random variable in the form
\[
    \vm{X} = \begin{pmatrix}
        \vm{Y} \\ \vmb{\nu}
    \end{pmatrix}
\]
where \(\vm{Y} \sim \Normal\para{\vmb{\lambda}, \vm{U}}\) and \(\vm{Y} \in \RR^m\). Hence, it has density
\[
    f_{\vm{Y}} \para{\vm{y}} = \frac{1}{\sqrt{\para{2 \pi}^m \det \vm{U}}} \exp \para{- \frac{\para{\vm{y} - \vmb{\lambda}}\transpose \vm{U}\inverse \para{\vm{y} - \vmb{\lambda}}}{2}}
\]
for the first \(m\) components.

For a Gaussian random variable, we have the converse of \autoref{claim:indep-diagonal-variance} is in fact true!

\begin{lemma}[Diagonal Variance implies Independence for Gaussian]
    \label{lem:diagonal-variance-independence-gaussian}%
    Let \(\vm{X} \sim \Normal\para{\vmb{\mu}, \vm{V}}\) where \(\vm{V}\) is positive definite. If \(\vm{V}\) is diagonal, then the \(X_i\)s are independent.
\end{lemma}

\begin{proof}[using Probability Density Function]
    Let
    \[
        \vm{V} = \diag\para{\lambda_1, \ldots, \lambda_n},
    \]
    and
    \[
        \mu = \para{\mu_1, \ldots, \mu_n}\transpose.
    \]

    The probability density function of \(\vm{X}\) is given by
    \begin{align*}
        f_{\vm{X}} \para{\vm{x}} &= \frac{1}{\sqrt{\para{2 \pi}^n \det \vm{V}}} \exp \para{- \frac{\para{\vm{x} - \vmb{\mu}} \transpose \vm{V}\inverse \para{\vm{x} - \vmb{\mu}}}{2}} \\
        &= \frac{1}{\sqrt{\para{2\pi}^n \det \vm{V}}} \cdot \exp\para{- \sum_{i = 1}^{n} \frac{\para{x_i - \mu_i}^2}{2 \lambda_i}},
    \end{align*}
    and the density function factorises. Therefore, \(X_i \sim \Normal\para{\mu_i, \lambda_i}\) is independent.
\end{proof}

\begin{proof}[using Moment Generating Function]
    We inspect the moment generation function. Recall that
    \[
        \vmb{\theta}\transpose \vm{X} \sim \Normal\para{\vmb{\theta}\transpose \vmb{\mu}, \vmb{\theta}\transpose \vm{V} \vmb{\theta}}.
    \]
    
    We have
    \begin{align*}
        m\para{\vmb{\theta}} &= \exp\para{\vmb{\theta}\transpose \vmb{\mu} + \frac{\vmb{\theta}\transpose \vm{V} \vmb{\theta}}{2}} \\
        &= \exp\para{\sum_{i = 1}^{n} \theta_i \mu_i + \sum_{i = 1}^{n} \frac{\theta_i^2 \lambda_i}{2}} \\
        &= \prod_{i = 1}^{n} \exp\para{\theta_i \mu_i + \frac{\theta_i^2 \lambda_i}{2}}
    \end{align*}
    factorised into moment generating functions for \(\Normal\para{\mu_i, \lambda_i}\). Therefore, \(X_1, \ldots, X_n\) are independent.
\end{proof}

Therefore, combining \autoref{claim:indep-diagonal-variance} and \autoref{lem:diagonal-variance-independence-gaussian}, we have that the components of a Gaussian vector is independent, if and only if their pairwise covariances are zero.

\subsection{Bivariate Gaussian Vector}

\begin{definition}[Bivariate Gaussian Vector]
    \label{def:bivariate-gaussian-vector}%
    The 2-Dimensional case of a multivariate Gaussian random variable, defined in \autoref{def:multivariate-gaussian-rv}, is known as a \emph{Bivariate Gaussian Vector}.
\end{definition}

\begin{proposition}[Conditional Expectation for Bivariate Gaussian]
    Let \(\para{X_1, X_2}\) be a bivariate Gaussian. Then we have
    \[
        \Expt\brac{\Conditional{X_2}{X_1}} = \mu_2 - a \mu_1 + a X_1,
    \]
    where
    \[
        a = \frac{\Cov\para{X_1, X_2}}{\Var\para{X_1}}.
    \]
\end{proposition}


\begin{proof}
    Let \(Y = X_2 - a X_1\) for some \(a \in \RR\). Then the vector
    \[
        \begin{pmatrix}
            X_1 \\ Y
        \end{pmatrix}
        =
        \begin{pmatrix}
            1 & 0 \\
            -a & 1
        \end{pmatrix}
        \begin{pmatrix}
            X_1 \\ X_2
        \end{pmatrix}
    \]
    is Gaussian. The covariance satisfies
    \[
        \Cov\para{X_1, Y} = \Cov\para{X_1, X_2 - a X_1} = \Cov\para{X_1, x_2} - a \Var\para{X_1}.
    \]
    
    Taking
    \[
        a = \frac{\Cov\para{X_1, X_2}}{\Var\para{X_1}}
    \]
    gives \(\Cov\para{X_1, Y} = 0\), which in turn makes them independent. Hence,
    \begin{align*}
        \Expt\brac{\Conditional{X_2}{X_1}} &= \Expt\brac{\Conditional{X_2 - a X_1}{X_1}} + \Expt\brac{\Conditional{a X_1}{X_1}} \\
        &= \Expt\brac{X_2 - a X_1} + a X_1 \\
        &= \mu_2 - a \mu_1 + a X_1.
    \end{align*}
\end{proof}